\section{Method}

\subsection{Conditions}

Four conditions vary only the collapse-prevention mechanism.

\begin{table}[h]
\centering
\small
\begin{tabular}{l c c l}
\toprule
\textbf{Condition} & \textbf{EMA target} & \textbf{Stop-gradient} & \textbf{Role} \\
\midrule
\texttt{ema\_stopgrad} & yes & yes & Classic JEPA baseline \\
\texttt{sigreg\_stopgrad} & no & yes & Both mechanisms \\
\texttt{sigreg\_nostopgrad} & no & no & LeJEPA's recipe; the condition under test \\
\texttt{none\_nostopgrad} & no & no & Sanity control; expected to collapse \\
\bottomrule
\end{tabular}
\end{table}

The \texttt{none\_nostopgrad} control removes every collapse-prevention mechanism at once:
one shared encoder, no detachment, no regularizer. Nothing prevents the constant solution,
so it must collapse. Its purpose is to demonstrate that the apparatus detects collapse when
collapse is present. If it fails to collapse, results from the other three conditions
cannot be trusted and the study is inconclusive.

An earlier version of this design paired an EMA target encoder with no stop-gradient as the
control. That combination is inert: with the EMA target's parameters held outside autograd,
the target branch reaches no trainable parameter, so disabling stop-gradient changes nothing
and the condition silently duplicates \texttt{ema\_stopgrad}. Removing collapse prevention
requires removing the teacher--student split as well.

\subsection{Shared training loop}

All four conditions execute the same training loop. The only swappable component is a
strategy object exposing three things: whether a separate EMA target encoder is
constructed, whether the target branch is detached, and how the total loss is assembled.

This is a correctness requirement rather than a code-organization preference. Had each
condition carried its own loop, a difference in outcomes could originate in an incidental
difference between loops rather than in the mechanism under study.

\subsection{Where the regularizer is applied}

SIGReg is applied to the pooled context-encoder embeddings, which always carry gradient.
Applying it to the target branch would make it a no-op in the stop-gradient conditions,
since a detached tensor propagates nothing. This choice is ours and is a deviation to note
when comparing against LeJEPA.
